% Soluzioni degli esercizi del Modulo 3 (generato da separa_soluzioni.py)
\chapter{Soluzioni -- Modulo 3 -- Array, Puntatori, Reference, Stack e Heap}
\textbf{Esercizio 3.1} --- Gli indirizzi cambiano a ogni esecuzione: il sistema
operativo carica il programma ogni volta in una posizione diversa.
\codice{codice/m03/es01.cpp}
\uscita{codice/m03/output/es01.txt}

\textbf{Esercizio 3.2}
\codice{codice/m03/es02.cpp}
\uscita{codice/m03/output/es02.txt}

\textbf{Esercizio 3.3} --- \cc|pm - v| vale 5: la differenza tra due puntatori
nello stesso array è il numero di \emph{elementi} che li separa.
\codice{codice/m03/es03.cpp}
\uscita{codice/m03/output/es03.txt}

\textbf{Esercizio 3.4}
\codice{codice/m03/es04.cpp}
\uscita{codice/m03/output/es04.txt}

\textbf{Esercizio 3.5}
\codice{codice/m03/es05.cpp}
\uscita{codice/m03/output/es05.txt}

\textbf{Esercizio 3.6}
\codice{codice/m03/es06.cpp}
\ingresso{codice/m03/input/es06.txt}
\uscita{codice/m03/output/es06.txt}

\textbf{Esercizio 3.7} --- \cc|contiene(v, i, v[i])| cerca \cc|v[i]| solo tra gli
elementi \emph{precedenti}: se c'è già, è un duplicato.
\codice{codice/m03/es07.cpp}
\uscita{codice/m03/output/es07.txt}

\textbf{Esercizio 3.8}
\codice{codice/m03/es08.cpp}
\uscita{codice/m03/output/es08.txt}

\textbf{Esercizio 3.9}
\codice{codice/m03/es09.cpp}
\uscita{codice/m03/output/es09.txt}

\textbf{Esercizio 3.10} --- Per inserire si sposta partendo dal \textbf{fondo}
(\cc|i = n| scendendo), altrimenti ogni elemento sovrascriverebbe il successivo
prima di averlo spostato. Per rimuovere si parte invece dalla posizione tolta.
\codice{codice/m03/es10.cpp}
\uscita{codice/m03/output/es10.txt}
